\subsubsection{分布的收敛}\label{subsec:distribution-convergence}

正如经典函数序列在偏微分方程中处于核心地位, 广义函数序列也是如此.

\begin{Definition}
\label{def:distribution-convergence}
若对每个 $\phi\in\mathcal{D}(\Omega)$ 都有 \[(f_n,\phi)\to(f,\phi)\] 则称 $\mathcal{D}^{\prime}(\Omega)$ 中的序列 $f_n$ \textbf{收敛}到 $f\in\mathcal{D}^{\prime}(\Omega)$.
\end{Definition}

\begin{Example}
\label{ex:uniform-convergence-distribution}
连续函数的一致收敛序列 (按 Example~\ref{ex:regular-distribution} 定义分布) 在 $\mathcal{D}^{\prime}$ 中也收敛.
\end{Example}

\begin{Example}
\label{ex:delta-rectangular-approximation}
考虑序列
\begin{equation}\label{eq:delta-approximating-rectangular-sequence}
f_n(x)=\begin{cases}n,&0<x<1/n\\0,&\text{otherwise.}\end{cases}
\end{equation}
有
\begin{equation}\label{eq:delta-rectangular-sequence-pairing}
\int_{-\infty}^{\infty}f_n(x)\phi(x)\,\mathrm{d}x=n\int_0^{1/n}\phi(x)\,\mathrm{d}x,
\end{equation}
当 $n\to\infty$ 时, 它收敛到 $\phi(0)$. 因而在 $\mathcal{D}^{\prime}(\mathbb{R})$ 中 $f_n(x)\to\delta(x)$.
\end{Example}

\begin{Remark}
\label{rem:distribution-density}
Problem~\ref{ex:compactly-supported-distribution-density} 要求读者证明每个分布都是紧支集分布的极限. 后面实际上将看到, 每个分布都是试验函数的极限; 换言之, 试验函数在 $\mathcal{D}^{\prime}(\Omega)$ 中稠密.
\end{Remark}

另一个重要结果是 $\mathcal{D}^{\prime}(\Omega)$ 的 (序列) 完备性.

\begin{Theorem}
\label{thm:distribution-sequential-completeness}
设 $f_n$ 为 $\mathcal{D}^{\prime}(\Omega)$ 中的序列, 使对每个 $\phi\in\mathcal{D}(\Omega)$, $(f_n,\phi)$ 都收敛. 那么存在 $f\in\mathcal{D}^{\prime}(\Omega)$, 使 $f_n\to f$.
\end{Theorem}

\begin{Proof}
定义
\begin{equation}\label{eq:distribution-pointwise-pairing-limit}
(f,\phi)=\lim_{n\to\infty}(f_n,\phi).
\end{equation}
显然, $f$ 是从 $\mathcal{D}(\Omega)$ 到 $\mathbb{R}$ 的线性映射. 为验证 $f\in\mathcal{D}^{\prime}(\Omega)$, 必须证明其连续性, 即必须证明若在 $\mathcal{D}(\Omega)$ 中 $\phi_n\to0$, 则 $(f,\phi_n)\to0$. 反设不成立. 那么, 选取一个子列, 仍将其记为 $\phi_n$,\footnote{对一个序列及其任意子列使用相同符号, 是偏微分方程中的典型做法. 其主要目的是使记号清楚 (因为经常必须考虑嵌套数层的子列), 但它还有一个令人愉快的副作用, 就是使许多经典分析学家发狂. 当然, 有时区分序列与其子列很重要 (如本证明后面所做的), 此时会使用适当的记号加以区分.} 可假设 $\phi_n\to0$, 但 $|(f,\phi_n)|\geq c>0$.

现在回顾, 在 $\mathcal{D}(\Omega)$ 中收敛到 0, 意味着所有 $\phi_n$ 的支集均位于 $\Omega$ 的某个固定紧子集内, 且 $\phi_n$ 的所有导数一致收敛到零. 再次选取子列后, 可假设对 $|\alpha|\leq n$, 有 $|D^\alpha\phi_n(\bm{x})|\leq4^{-n}$. 现在令 $\psi_n:=2^n\phi_n$. 那么 $\psi_n$ 仍在 $\mathcal{D}(\Omega)$ 中收敛到 0, 但 $|(f,\psi_n)|\to\infty$.

现在递归构造 $\{f_n\}$ 的子列 $\{f_n'\}$ 及 $\{\psi_n\}$ 的子列 $\{\psi_n'\}$. 首先选取 $\psi_1'$, 使 $|(f,\psi_1')|>1$. 由于 $(f_n,\psi_1')\to(f,\psi_1')$, 可选取 $f_1'$, 使 $|(f_1',\psi_1')|>1$. 现在假设已对 $j<n$ 选取了 $f_j'$ 和 $\psi_j'$. 然后从序列 $\{\psi_n\}$ 中选取 $\psi_n'$, 使
\begin{equation}\label{eq:distribution-completeness-small-earlier-pairings}
|(f_j',\psi_n')|<\dfrac1{2^{n-j}},\quad j=1,2,\ldots,n-1,
\end{equation}

\begin{equation}\label{eq:distribution-completeness-large-limit-pairing}
|(f,\psi_n')|>\sum_{j=1}^{n-1}|(f,\psi_j')|+n.
\end{equation}
这是可能的, 因为一方面 $\psi_n\to0$, 另一方面 $|(f,\psi_n)|\to\infty$. 此外, 由于 $(f_n,\psi)\to(f,\psi)$, 可选取 $f_n'$, 使
\begin{equation}\label{eq:distribution-completeness-diagonal-pairing}
|(f_n',\psi_n')|>\sum_{j=1}^{n-1}|(f_n',\psi_j')|+n.
\end{equation}
接着令
\begin{equation}\label{eq:distribution-completeness-test-series}
\psi=\sum_{n=1}^{\infty}\psi_n'.
\end{equation}
由 $\psi_n'$ 的构造可知, 右端级数在 $\mathcal{D}(\Omega)$ 中收敛. 因而
\begin{equation}\label{eq:distribution-completeness-pairing-series-split}
(f_n',\psi)=\sum_{j=1}^{n-1}(f_n',\psi_j')+(f_n',\psi_n')+\sum_{j=n+1}^{\infty}(f_n',\psi_j').
\end{equation}
由 \eqref{eq:distribution-completeness-small-earlier-pairings} 得到
\begin{equation}\label{eq:distribution-completeness-tail-estimate}
\left|\sum_{j=n+1}^{\infty}(f_n',\psi_j')\right|<\sum_{j=n+1}^{\infty}2^{n-j}=1,
\end{equation}
将这一点与 \eqref{eq:distribution-completeness-pairing-series-split} 和 \eqref{eq:distribution-completeness-diagonal-pairing} 结合, 得出 $|(f_n',\psi)|>n-1$. 因而当 $n\to\infty$ 时, $(f_n',\psi)$ 的极限不存在, 矛盾.
\end{Proof}

可以用类似的反证法证明以下引理; 证明细节留作习题 (参见 Problem~\ref{ex:distribution-test-sequence-pairing}).

\begin{Lemma}
\label{lem:distribution-test-sequence-pairing}
假设在 $\mathcal{D}^{\prime}(\Omega)$ 中 $f_n\to0$, 且在 $\mathcal{D}(\Omega)$ 中 $\phi_n\to0$. 那么 $(f_n,\phi_n)\to0$.
\end{Lemma}

还有以下推论.

\begin{Corollary}
\label{cor:distribution-joint-continuity}
若 $f_n\to f$ 且 $\phi_n\to\phi$, 则 $(f_n,\phi_n)\to(f,\phi)$.
\end{Corollary}

因而分布与试验函数之间的配对连续. (当然, 由定义可明显看出对每个因子的分别连续性, 但联合连续性需要证明.)

\begin{Proof}
推论直接由恒等式得出
\begin{equation}\label{eq:distribution-joint-pairing-identity}
(f_n-f,\phi_n-\phi)=(f_n,\phi_n)-(f,\phi_n)-(f_n,\phi)+(f,\phi).
\end{equation}

\end{Proof}
